\section{Research Methodology}\label{sec:methodology}

This research follows a controlled experimental methodology to evaluate
the contribution of different components in a hybrid stock market
forecasting system. Four forecasting models are developed, where each
successive model differs from the previous one by only a single
architectural modification. All models are tested using the same
dataset, identical prediction dates, and the same evaluation samples to
ensure a fair comparison. This section describes the dataset, data
splitting strategy, forecasting models, retrieval process, and
evaluation metrics used throughout the study.

\subsection{Dataset}\label{sec:dataset}

Historical daily market data for five U.S. stocks---\textbf{NVIDIA
(NVDA), Apple (AAPL), Tesla (TSLA), JPMorgan Chase (JPM), and Microsoft
(MSFT)}---were collected using the \textbf{yfinance} library. The
dataset covers the period from \textbf{1 January 2019 to 31 December
2024}, providing \textbf{1,509 trading days} for each company. Daily
values of the \textbf{CBOE Volatility Index (VIX)} were also collected
over the same period to represent overall market conditions.

The selected companies represent different market behaviours rather than
a single market segment. For example, NVIDIA experienced exceptionally
high returns and volatility during the study period, whereas JPMorgan
showed comparatively lower growth and more stable price movements.
Including stocks with different volatility characteristics allows the
proposed framework to be evaluated under diverse market conditions
rather than only during stable periods.

The prediction task is formulated as a \textbf{binary classification
problem}. For each trading day, the model predicts whether the stock
closing price will increase or decrease after \textbf{five trading
days}. A value of \textbf{1} represents an increase in price, while
\textbf{0} represents a decrease. Predicting market direction is
considered more suitable for trading decisions than forecasting the
exact future price because stock returns are generally centred around
zero and small price changes can dominate regression-based evaluation.

To represent market behaviour, \textbf{22 technical features} are
generated for each stock on every trading day. These include the
Relative Strength Index (RSI), Moving Average Convergence Divergence
(MACD) and its signal line, moving-average ratios, Bollinger Band
indicators, historical returns over different time windows, short-term
volatility, volume ratio, Average True Range (ATR), VIX-related
indicators, and features describing the relationship between each stock
and an equally weighted portfolio of all five stocks. All features are
calculated exclusively from historical information available on or
before the prediction date, ensuring that no future information is
introduced into the model.

\begin{quote}
\textbf{Market Regime Classification}
\end{quote}

To analyse model performance under different market conditions, each
trading day is assigned to one of two market regimes based on the daily
VIX value. Days with a \textbf{VIX value greater than 25} are classified
as \textbf{SHOCK} periods, while all remaining days are classified as
\textbf{CALM} periods. Since the VIX value is known on the same trading
day, this classification introduces no future information. Model
performance is reported separately for both regimes because prediction
difficulty and baseline market behaviour differ significantly between
stable and highly volatile periods.

\subsection{Data Splitting Protocol}\label{sec:data-splitting}

The LSTM forecasting model is trained using a \textbf{walk-forward
validation} strategy consisting of three chronological folds. For every
testing period, only historical observations occurring before the
prediction date are used for model training. A \textbf{five-day embargo
period} is inserted between the training and testing datasets to prevent
overlap between training labels and future test labels, thereby
eliminating potential information leakage.

Within each training fold, the most recent \textbf{15\%} of the training
data is reserved as a validation set. This validation data is used for
early stopping and decision-threshold selection. Random shuffling is
intentionally avoided because it would violate the temporal ordering of
financial data and could introduce unrealistic validation results.
Feature normalization is performed using parameters computed only from
the training data, and these parameters are subsequently applied to the
validation and testing datasets.

For the comparative evaluation of the four forecasting systems,
\textbf{60 trading dates} are selected from the walk-forward testing
period. Predictions are generated for all five stocks on each selected
date, resulting in \textbf{300 independent evaluation samples},
including \textbf{190 calm-market cases} and \textbf{110 shock-market
cases}. An additional tuning dataset containing \textbf{200 separate
samples} is reserved for selecting model thresholds and fusion
parameters. This tuning dataset is completely independent of the final
evaluation dataset to ensure that the reported experimental results
remain unbiased.

\subsection{Forecasting Models}\label{sec:forecasting-models}

The proposed framework evaluates four forecasting approaches under
identical experimental conditions.

\begin{itemize}
\item
  \textbf{Model 1 (ML):} A two-layer LSTM network processes the previous
  \textbf{60 trading days} of the 22 engineered features and outputs the
  probability of an upward market movement.
\item
  \textbf{Model 2 (LLM):} Gemini 2.5 Flash is provided with the same
  market information converted into structured textual format. The
  prompt includes normalized closing prices, RSI, short- and medium-term
  returns, volatility, trading volume ratio, and MACD ratio. The model
  is required to predict market direction, provide a confidence score,
  generate a brief explanation, and report the numerical values used
  during reasoning.
\item
  \textbf{Model 3 (LLM + RAG):} This model uses the same prompt as the
  standalone LLM but additionally receives relevant historical evidence
  retrieved from an external knowledge base before generating its
  prediction.
\item
  \textbf{Model 4 (Hybrid):} The final model combines the prediction
  probabilities produced by the LSTM and the grounded language model
  through a calibrated fusion strategy.
\end{itemize}

\begin{quote}
\textbf{Stock Anonymization}
\end{quote}

To minimise the possibility of the language model relying on memorized
historical knowledge, company names are replaced with anonymous
identifiers such as \textbf{"Stock A"} during inference. Since Gemini
has been trained on publicly available information covering the
evaluation period, this anonymization helps ensure that predictions are
based on the provided market data rather than previously memorized
events. The effectiveness of this strategy is verified experimentally
and discussed in the evaluation section.

\begin{quote}
\textbf{Evidence Retrieval}
\end{quote}

The retrieval system uses a manually verifiable knowledge corpus
containing \textbf{252 evidence entries}. The corpus consists of three
categories: documented macroeconomic events, automatically generated
market stress events triggered by high VIX values, and automatically
detected large movements in the equally weighted stock portfolio. Except
for the manually collected macroeconomic events, all remaining evidence
is generated directly from historical market data.

Relevant evidence is retrieved using \textbf{TF-IDF cosine similarity}
combined with a time-based weighting function that gradually reduces the
importance of older documents within a \textbf{90-day} historical
window. The four most relevant evidence items are provided to the
language model for each prediction. To eliminate look-ahead bias, only
documents published on or before the prediction date are considered
during retrieval.

\subsection{Evaluation Metrics}\label{sec:evaluation-metrics}

The proposed forecasting framework is evaluated using five complementary
experiments.

\begin{quote}
\textbf{Prediction Accuracy and Statistical Significance}
\end{quote}

Prediction accuracy is estimated using \textbf{2,000 bootstrap
resamples} generated at the date level rather than the individual
prediction level because multiple stock predictions made on the same day
are statistically correlated. Confidence intervals are calculated using
these bootstrap samples, while pairwise comparisons between forecasting
models are performed using the \textbf{exact McNemar test}, which is
appropriate for paired classification results. Probability calibration
is evaluated using the \textbf{Brier Score}.

\begin{quote}
\textbf{Hallucination Analysis}
\end{quote}

The reliability of language model outputs is evaluated through an
automated hallucination assessment. The numerical values reported by the
language model are compared directly with the input values supplied in
the prompt. Incorrect numerical values are recorded, while scale-related
errors are classified separately as unit errors. The evaluation also
identifies invented evidence citations and checks whether the language
model reveals the actual company names despite anonymization. The entire
evaluation process is fully automated and does not require human
judgement or another language model.

\begin{quote}
\textbf{Prediction Stability}
\end{quote}

Model consistency is measured by submitting the same prompt \textbf{five
times} for each of \textbf{30 evaluation cases}. Differences between
repeated predictions are recorded to quantify output variability. This
experiment evaluates the stability of each forecasting system rather
than assuming deterministic behaviour simply because the language model
operates with a temperature value of zero.

\begin{quote}
\textbf{Response Quality Assessment}
\end{quote}

The usefulness of each generated response is evaluated using a
\textbf{17-criterion scoring framework} covering four dimensions:
\textbf{grounding (30 points), personalization (25 points), helpfulness
(25 points), and safety (20 points)}. The evaluation is performed using
deterministic rule-based scoring instead of language-model-based judges
to ensure complete reproducibility. Only text generated by the
forecasting models is evaluated, while fixed template text used by the
hybrid system is excluded from scoring.

\begin{quote}
\textbf{Ceiling Analysis}
\end{quote}

Finally, an upper-bound analysis is performed to estimate the maximum
achievable performance for the forecasting task. Four performance bounds
are computed using the same evaluation dataset, including threshold
optimization on the testing data, optimized hybrid combinations, and an
ideal oracle that always selects the correct prediction among the four
models. Although the oracle cannot be implemented in practice because it
requires knowledge of the true outcome, it provides a theoretical upper
limit for evaluating the effectiveness of the proposed forecasting
approaches.
